% !TeX root = ../main.tex
% sections/02_related.tex -- Section 2 of the MobiWac 2026 camera-ready (accepted 2026-08-26).
% Base text: the dissertation's Chapter 5 (articles/dissertacao/src/chapters/5_mobiwac/02_related.tex),
% ported 2026-09-06 (CONSOLIDATION_PLAN 3.11). The chapter's commented-out recap subsection
% (dissertation-only bridging prose, never rendered in either text) was dropped from this copy;
% the chapter keeps it. Edits made for the camera-ready carry a dated comment beside them.
\section{Background and Related Work}
\label{sec:related}

\subsection{From place embeddings to per-visit embeddings}
\label{sec:related-embeddings}
A place embedding transforms a POI into a vector so that similar or nearby
places sit close together in the vector space.
Deep Graph Infomax (DGI)~\cite{velickovic2019dgi}, a general self-supervised
method for graphs, learns such vectors on a place network linked by
similarity, time, and distance.
Hierarchical Graph Infomax (HGI) \cite{huang2023hgi} adds geographic structure,
organizing the network as place, then region, then city. Both produce one vector per place, so two visits to the same coffee shop look
identical to the model.
Contextual check-in representations instead give each visit its own vector: CTLE
\cite{lin2021ctle}, the closest example, learns a vector per visit by masking and
reconstructing parts of a user's check-in sequence. CTLE is a sequence model, a Transformer that reads the
check-in sequence itself; our representation, Check2HGI, remains a network model, keeping the
hierarchical place-to-region-to-city network and the same infomax objective,
now extended one level deeper, from individual places down to individual
check-ins (Fig.~\ref{fig:dataflow}). The order of a user's visits still enters
the network, as links between consecutive check-ins
(Section~\ref{sec:method-rep}).

We compare against CTLE and a feature-concatenation control to separate the
combination from contextualization alone and from feature injection.

\subsection{Predicting the next category and the next region}
\label{sec:related-tasks}
A large body of work predicts the next place that a user will visit, the exact POI,
with recurrent and attention models such as DeepMove \cite{feng2018deepmove}, STAN
\cite{luo2021stan}, and GETNext \cite{yang2022getnext}. In contrast, we target two properties of the next visit, its category and its
region, rather than the exact place. We choose them because they are
what a mobility-aware service can act on, not to make the task easier; the region
alone spans hundreds to several thousand classes. Predicting
over a partition of the map is standard in the human-mobility literature, with a grid cell as the
target~\cite{luca2021mobilitysurvey}; we substitute official neighborhood-scale units. Our earlier
work~\cite{silva2025mtlnet} paired static category classification with next-category prediction. This paper introduces the next-region task, adds the check-in-level
representation, and reports, for each task, where the joint model outperforms a dedicated single-task model
and where their difference stays within a stated bound (Section~\ref{sec:results-part2}).
% [2026-09-06, camera-ready, S9] REWRITTEN, not deleted: GLOSSARY section 9.3 allows exactly one
% self-positioning sentence in this section. The old clause "on which sharing helps instead of
% hurting" was the v17 verdict map (category gain at every dataset). Under the reported results
% (CAMERA_READY section 3) sharing outperforms the dedicated model in three of twelve comparisons
% (region at Texas and California, category at Florida); the remaining differences are bounded
% (region: the registered two-point margin; category: the derived half-point bound). The new
% clause states the delta over the earlier work without a verdict, and points at the section
% that carries the verdicts. Same sentence in both trees before this edit: the dissertation
% chapter keeps the old wording (it is delivered text); this divergence is declared in ERRATA.md.

The field
increasingly models several granularities at once; in those systems, category and
region are auxiliary signals that help a primary next-place task
(MCMG \cite{sun2024mcmg}, HMT-GRN \cite{lim2022hmtgrn}).
We study the pair as the object itself. To our knowledge, fine-grained region as an
end target of equal standing, rather than an auxiliary coarse grid cell, is
underexplored.
The nearest exceptions do not study our exact pairing:
DRRGNN~\cite{zhu2022drrgnn} forecasts a person's next activity region jointly
with a mobility-intention label, but over regions discovered per person rather
than a fixed citywide partition; a generative recommender~\cite{sun2025kgtb}
predicts category and region together, but only as auxiliary steps toward its
next-place ranking.

\subsection{Multi-task learning for mobility}
\label{sec:related-mtl}
MCARNN \cite{Liao2018} jointly predicts activity and
location, and a recent hierarchy-aware model continues the
pattern \cite{wang2025hamtl}; in both, the location target is
the exact place. The closest alternative to our setting is
the cascade: predict the category first, then the place given the
category \cite{ye2013nextmove}. CSLSL
\cite{huang2024cslsl} is its current form, predicting in a chain (when, then what,
then where), with location as the primary output and category as an
intermediate step; CatDM \cite{yu2020catdm} likewise uses the predicted
category only to filter candidate places. We instead predict
category and region in parallel, as end targets of equal standing in one model
with a single forward pass, and we drop the next-place target entirely. On optimization, we are conservative by design: joint training with a fixed loss weighting is
standard practice~\cite{caruana1997multitask}, and gradient-balancing methods (PCGrad~\cite{yu2020pcgrad}, Nash-MTL~\cite{navon2022nashmtl})
rarely improve on a well-tuned fixed weighting with two tasks~\cite{xin2022domtl}. We confirmed
this during development, on an earlier preparation of the data: a screen of nineteen loss and
gradient balancers at their default configurations, at a single seed on Alabama and Florida, found
none that improved on fixed equal weighting on both tasks at both datasets.  The cosine similarity between the next-category and next-region updates on the
shared trunk, measured on the reported joint model at four of the six datasets (Istanbul, Alabama,
Arizona, and Florida), is equivalent to zero at every one of them, with per-dataset means within
two thousandths of zero, so a balancer has no conflict to resolve. This is a finding for this pair
of tasks, not a general rule.
% [2026-09-06, cut C5] Two paragraphs (~330 words) -> one (~150). Kept, as the mobiwac list
% requires: PCGrad and Nash-MTL by name, "on an earlier preparation of the data" for the screen,
% and the cosine finding on the reported model (four of the six datasets named; "within two
% thousandths of zero", apx_f_cosine.tex:339-340). Dropped: the two named exceptions of the screen
% (Nash-MTL and scale normalization at Alabama and their Florida outcomes; T4_full_screen.json,
% pre-v18 values, kept in the provenance comment history of this file up to 2026-09-06) and the
% development-time four-Gowalla-state cosine (+0.001 / +0.0032, incl. Georgia), superseded by the
% reported-model measurement. No claim strengthened: the screen's verdict is still "none improved
% on both tasks at both datasets"; the cosine is still "equivalent to zero".
% [2026-09-06, camera-ready] "four datasets" -> "four of the six datasets (Istanbul, Alabama,
% Arizona, and Florida)": the previous sentence's "four Gowalla states" is a different set
% (Alabama, Arizona, Florida, Georgia), so an unqualified second "four" read as the same pool.
% Names from articles/dissertacao/src/chapters/apx_f_cosine.tex:191-192 ("four datasets: Istanbul,
% Alabama, Arizona, and Florida"); "within two thousandths of zero" from :339-340. No number changed.
% [round9f, 2026-08-02, PENDENCIAS_RESOLVIDOS 2.9 (arquivado 2026-08-02), AUTHOR DECISION]
% ADDED HERE AND IN THE [mobiwac] ARTICLE, in the same words, so the two texts stay
% identical as the round-6 note below requires. The author's instruction was to write the
% sentence so it SURVIVES STANDALONE -- no reference to a document the reader may not hold --
% because this article has no appendix and a cross-document pointer cannot be resolved from
% here. The dissertation carries the pointer to its own Appendix D in the paragraph that
% introduces the section, not in this sentence.
% NUMBERS QUOTED, NOT COMPUTED: the +0.001 / +0.0032 figures above are the EARLIER four-seed
% development measure and are left exactly as they were. [2026-09-06] The sentence that once read
% "positive at every one of them" over a seven-dataset pool was replaced on 2026-08-20 (65ada3cd) by the
% four-dataset equivalence result; the old wording is on the never-cite list (CAMERA_READY section 4).
% [v2 review F-03, 2026-07-26] SCOPE CORRECTION (persona 10). The sentence said "three of our six
% datasets". The measurement pools FOUR Gowalla states -- alabama, arizona, georgia, florida --
% per docs/results/mtl_improvement/R0_matched_metric_bar.json (its "states" object has exactly those
% four keys, each with the four seeds 0/1/7/100) and scripts/.../plot_grad_cosine.py:19-24, whose
% STATE_STYLE map names the same four. Georgia is NOT one of the six datasets this dissertation
% reports, so "three of our six" both undercounted the pool and implied every state in it is one we
% report. Read from the JSON, not recomputed. The correction is in the author's favour: the finding
% replicates on a state the dissertation never reports.
% [round6, 2026-07-28] PARITY CLOSED. That F-03 correction had been applied on the dissertation side
% only, so until now the submitted paper still read "three of our six datasets" two lines from the
% sentence above, while this paper named the four states. Under the Chapter 5 regime the two texts
% must be identical, so the dissertation's wording was carried into
% articles/[mobiwac]/src/sections/02_related.tex this round, with "this dissertation" rendered as
% "this study" there. WHICH TEXT WAS RIGHT: this one. I re-read the source myself rather than trusting
% the earlier correction. R0_matched_metric_bar.json has exactly four keys under "states" -- alabama,
% arizona, georgia, florida -- and "seeds": [0, 1, 7, 100]; scripts/mtl_improvement/plot_grad_cosine.py
% reads its per-state run directories out of that same JSON and its STATE_STYLE map names the same
% four. So four Gowalla states at four seeds each, of which three are among the five United States
% datasets this document reports. "Three of our six datasets" and "four Gowalla states" are different
% claims about the same measurement, and only the second one matches the artifact. Logged in
% articles/[mobiwac]/ERRATA.md.
% [round5, author decision 2.4] The text read as a two-optimizer test when the sweep was much
% wider. Source of record: docs/results/mtl_improvement/T4_full_screen.json, 19 methods x 2 states:
%   uncertainty_weighting, equal_weight, static_weight, scale_norm, scheduled_static, uw_so, dwa,
%   famo, gradnorm, db_mtl, cagrad, nash_mtl, aligned_mtl, pcgrad, fairgrad, stch, go4align,
%   bayesagg_mtl, excess_mtl.
% AL next-category macro-F1 vs equal_weight 53.57: TWO methods are above it, nash_mtl 54.25 (+0.68)
% and scale_norm 53.76 (+0.19). [correction] An earlier version of this line and of the sentence said
% nash_mtl was the ONLY one, while this very comment listed scale_norm as also above -- a claim that
% contradicted its own ledger. Caught in review. Full picture, all four cells:
%   nash_mtl:   AL cat 54.25 (+0.68) reg 62.77 (+0.24) | FL cat 71.61 (-0.19) reg 73.04 (-0.02)
%   scale_norm: AL cat 53.76 (+0.19) reg 61.96 (-0.57) | FL cat 72.12 (+0.33) reg 35.47 (-37.59)
% nash_mtl is the only one above on BOTH tasks at AL; scale_norm is the only one above on category
% at BOTH datasets, and its region score collapses. So no method wins everywhere, which is why
% the sentence says "across both tasks and both datasets" and then names both exceptions with
% what each one gives up. PCGrad 53.11 (-0.46) is kept by name at the
% author's request so an MTL reviewer sees it was tried.

